\documentclass[CJK]{beamer}
\usepackage{CJKutf8}
\usepackage{beamerthemesplit}
\usetheme{Malmoe}
\useoutertheme[footline=authortitle]{miniframes}
\usepackage{amsmath}
\usepackage{amssymb}
\usepackage{graphicx}
\usepackage{color}
\usepackage{slashed}
\usepackage{simplewick}
\graphicspath{{../figures/}}
\def\be{\begin{equation}}
\def\ee{\nonumber\end{equation}}
\def\bea{\begin{eqnarray}}
\def\eea{\nonumber\end{eqnarray}}
\def\ii{{\dot{\imath}}}
\def\bch{\begin{CJK}{UTF8}{gbsn}}
\def\ech{\end{CJK}}
\def\bex{\begin{minipage}{0.3\textwidth}\includegraphics[width=1in]{jugelizi.png}\end{minipage}\begin{minipage}{0.6\textwidth}}
\def\eex{\end{minipage}}
\def\chtitle#1{\frametitle{\bch#1\ech}}
\def\skipline{{\vskip0.1in}}
\def\skiplines{{\vskip0.2in}}
\def\lagr{{\mathcal{L}}}
\def\hamil{{\mathcal{H}}}
\def\vecv{{\mathbf{v}}}
\def\vecx{{\mathbf{x}}}
\def\veck{{\mathbf{k}}}
\def\vecp{{\mathbf{p}}}
\def\vecn{{\mathbf{n}}}
\def\vecA{{\mathbf{A}}}
\def\vecP{{\mathbf{P}}}
\def\vecsigma{{\mathbf{\sigma}}}
\def\hatJn{{\hat{J_\vecn}}}
\def\hatJx{{\hat{J_x}}}
\def\hatJy{{\hat{J_y}}}
\def\hatJz{{\hat{J_z}}}
\def\hatj#1{\hat{J_{#1}}}
\def\hatphi{{\hat{\phi}}}
\def\hatq{{\hat{q}}}
\def\hatpi{{\hat{\pi}}}
\def\vel{\upsilon}
\def\Dint{{\mathcal{D}}}
\def\adag{{\hat{a}^\dagger}}
\def\bdag{{\hat{b}^\dagger}}
\def\cdag{{\hat{c}^\dagger}}
\def\ddag{{\hat{d}^\dagger}}
\def\hata{{\hat{a}}}
\def\hatb{{\hat{b}}}
\def\hatc{{\hat{c}}}
\def\hatd{{\hat{d}}}
\def\hatN{{\hat{N}}}
\def\hatH{{\hat{H}}}
\def\hatp{{\hat{p}}}
\def\Fup{{F^{\mu\nu}}}
\def\Fdown{{F_{\mu\nu}}}
\def\newl{\nonumber \\}
\def\SIkm{\mathrm{km}}
\def\SIyr{\mathrm{yr}}
\def\SIGyr{\mathrm{Gyr}}
\def\SIeV{\mathrm{eV}}
\def\SIGeV{\mathrm{GeV}}
\def\SIm{\mathrm{m}}
\def\SIcm{\mathrm{cm}}
\def\SIJ{\mathrm{J}}
\def\SIs{\mathrm{s}}
\def\SIkg{\mathrm{kg}}
\def\SIg{\mathrm{g}}
\def\vece{\mathrm{e}}
\def\bmat#1{\left(\begin{array}{#1}}
\def\emat{\end{array}\right)}
\def\bcase#1{\left\{\begin{array}{#1}}
\def\ecase{\end{array}\right.}
\def\calM{{\mathcal{M}}}
\def\calT{{\mathcal{T}}}
\def\calR{{\mathcal{R}}}
\def\barpsi{\bar{\psi}}
\def\baru{\bar{u}}
\def\barv{\bar{\upsilon}}
\def\bmini#1{\begin{minipage}{#1\textwidth}}
\def\emini{\end{minipage}}
\def\qeq{\stackrel{?}{=}}
\def\torder#1{\mathcal{T}\left(#1\right)}
\def\rorder#1{\mathcal{R}\left(#1\right)}


\title{Quantum Field Theory I \\ Lesson 02 - Action Princple}
  \author{}
  \date{}


\begin{document}

\begin{frame}
 
\begin{center}
\begin{Large}
\bch
量子场论 I 

{\vskip 0.3in}

第二课 作用量原理和Noether定理

\ech
\end{Large}
\end{center}

\vskip 0.2in

\bch
课件下载
\ech
https://github.com/zqhuang/SYSU\_QFTI

\end{frame}


\begin{frame}
\frametitle{\bch 课堂互动：函数与稳定点 \ech}
\bch
稳定点：所有一阶导数全为零的点 （极值点一定是稳定点，反之则未必）
\begin{itemize}
\item{$f(\phi) = \phi$存在稳定点吗}
\item{求$f(\phi) = \phi^2$的稳定点}
\item{求$f(\phi) = \frac{1}{3}\phi^3-\phi$的稳定点}
\item{求$f(\phi_1, \phi_2) = \frac{1}{3}\phi_1^3 - \phi_1 + (\phi_2-\phi_1)^2$的稳定点}
\item{求$f(\phi_1, \phi_2, \ldots, \phi_n) = \frac{1}{3}\phi_1^3 - \phi_1 + \sum_{j=2}^n(\phi_j - \phi_1)^2$的稳定点}
\item{求$f(\phi_1, \phi_2, \ldots) = \frac{1}{3}\phi_1^3 - \phi_1 + \sum_{j=2}^\infty(\phi_j - \phi_1)^2$的满足“边界条件”$\phi_1>0$的稳定点。}
\end{itemize}

\skipline

我们看到函数如存在多个稳定点，则可以添加适当的“边界条件”使得稳定点唯一。

\ech
\end{frame}

\begin{frame}
\chtitle{场的函数}
\bch
以自然数标记自由度的变量$(\phi_1, \phi_2, \ldots)$的进一步推广就是用坐标来标记自由度的“场”$\phi_t$ ($t\in (-\infty, \infty)$)。当然，我们更习惯把它记成$\phi(t)$。在学习场论的过程中，我们通常不把场看成坐标的函数，而是把坐标$t$看成自由度的标记，把场看成有无穷多自由度的变量。

\skipline

场既然只是一个无穷多元变量，当然就可以定义它的函数，数学上叫“泛函”(functional)：
\ech
\bex
\bch
$$ S = \int_{-\infty}^\infty \frac{m}{2}\left[(\frac{d\phi}{dt})^2 - \omega^2\phi^2\right] dt$$
把$\phi(t)$映射为一个数。这里$\omega$是具有角速度量纲的常量，$m$是具有质量量纲的常量。
\ech
\eex
\end{frame}

\begin{frame}
\chtitle{四维时空里的场的泛函}
\bch
四维时空里的场（也同样可以看成一个无穷多自由度的变量，其自由度用四维时空坐标来标记）也可以有它的泛函，例如。
\ech
\bex
$$ S = \int \sqrt{-g}d^4x\, \left[\frac{1}{2}\partial^\mu\phi \partial_\mu\phi - \frac{1}{2}m^2\phi^2\right]$$
\bch
把场$\phi(x^0,x^1, x^2, x^3)$映射为一个数$S$. 其中$g$是度规$g_{\mu\nu}$的行列式的简写。$d^4x$是$dx^0dx^1dx^2dx^3$的简写。$m$是质量量纲的常量。积分范围是整个四维时空。
\ech
\eex
\end{frame}

\begin{frame}
\chtitle{作用量}
\bch
泛函

$$ S = \int_{-\infty}^\infty \frac{m}{2}\left[(\frac{d\phi}{dt})^2 - \omega^2\phi^2\right] dt$$

定义了一个物理系统(谐振子）。反过来的等价说法是这个物理系统的“作用量”（Action）是上述泛函。

\skipline
{\bf 作用量的满足边界条件的稳定点给出了物理系统的经典解。}原则上所有经典物理问题都可以归结为求解作用量的稳定点。
\ech

\end{frame}


\begin{frame}
\chtitle{课堂互动}
\bch
求谐振子
$$ S = \int_{-\infty}^\infty \frac{m}{2}\left[(\frac{d\phi}{dt})^2 - \omega^2\phi^2\right] dt$$
的运动方程。在边界条件$\left.\phi\right\vert_{t=0} = 1, \left.\frac{d\phi}{dt}\right\vert_{t=0} = 0$下写出它的经典解$\phi(t)$.
\ech

\end{frame}


\begin{frame}
\chtitle{课堂互动}
\bch
考虑一维空间x和一维时间t构成的时空里的标量场$\phi(x,t)$，若其作用量为
$$ S = \int dx dt \   \frac{1}{2} \left[\left(\frac{\partial \phi}{\partial t}\right)^2 - \left(\frac{\partial \phi}{\partial x}\right)^2 - m^2\phi^2\right]\, ,$$
求$\phi(x, t)$的运动方程。
\ech

\end{frame}


\begin{frame}
\chtitle{Euler-Lagrange方程}
\bch
一般性地，作用量为
$$S = \int d^4x\,\lagr(\phi,\partial_\mu\phi)$$
的运动方程（即Euler-Lagrange方程）为
$$\partial_\mu\left(\frac{\partial \lagr}{\partial\left(\partial_\mu \phi\right)} \right)-\frac{\partial \lagr}{\partial \phi} = 0$$
其中$\lagr$称为拉氏密度(Lagrangian density).注意拉氏密度在形式上必须写成场和场的时空偏导的函数。

\skipline
注意：推导Euler-Lagrange方程并不像很多教科书上那样需要对边界条件进行一些假设。
\ech
\end{frame}

\begin{frame}
\chtitle{时间维特殊化}
\bch
如果想回到我们习惯的“时间”的概念，可以把把时间维特殊化。把场$\phi$看成用三维坐标标记的无穷多元变量$\phi(\vecx)$随时间的演化。
把作用量写成拉氏量（Lagrangian, 定义为拉氏密度的空间积分$L \equiv \int d^3\vecx\, \lagr$）的时间积分。
$$S = \int L(\phi,\dot\phi) dt$$
这里$\dot\phi$是$\phi$的时间偏导数（即固定$\vecx$处的$\phi$的时间导数）。注意：一旦我们采用这种把时间维特殊化的描述方法, $\phi$的空间微分在形式上就只是看成不同空间点的$\phi$联合得到的函数，而不看成$L$所依赖的变量了。

在这种描述方法下，Euler-Lagrange方程简化为：
$$\frac{d}{dt} \frac{\partial L}{\partial \dot\phi(\vecx)} = \frac{\partial L}{\partial \phi(\vecx)}\,.$$

\skipline

\ech
\end{frame}


\begin{frame}
\chtitle{Hamilton量}
\bch
把时间维度$t$特殊化后，$\phi(\vecx)$共轭的正则动量定义为：
$$ \pi(\vecx) \equiv \left.\frac{\partial\lagr}{\partial\dot\phi}\right\vert_{\vecx} \,.$$
显然$\pi(\vecx)$也随时间变化。

场的Hamilton密度定义为
$$ \hamil \equiv \pi \dot\phi- \lagr $$
注意Hamilton密度形式上必须写成正则变量$\phi$（包括空间偏导都必须离散化后写成$\phi$的函数)和正则动量$\pi$的函数,时间导数$\dot\phi$不允许出现在$\hamil$的最终表达式里。

场的Hamilton函数$H$则是$\hamil$的空间积分。
$$ H(t) = \int d^3\vecx\, \hamil\left(\phi(t,\vecx), \pi(t,\vecx)\right)\,.$$
\ech
\end{frame}

\begin{frame}
\chtitle{Hamilton方程}
\bch
理论力学里的Hamilton方程对有无穷多自由度的场依然成立, 在任意时间$t$：
$$\dot\pi(\vecx) = -\frac{\partial \hamil}{\partial\phi(\vecx)},\ \dot\phi(\vecx) = \frac{\partial\hamil}{\partial\pi(\vecx)}$$.
\ech
\end{frame}

\begin{frame}
\chtitle{Noether定理}
\bch
仍回到四维时空的作用量。

$$S = \int d^4x\ \lagr(\phi, \partial_\mu\phi)\,.$$

考虑场的一个无穷小变化：

$$\delta \phi = \epsilon \frac{\delta\phi}{\delta\epsilon}$$

若这个映射使得作用量的变化严格为零，则$\lagr$的变化必然可以写成$ \delta \lagr = \epsilon \partial_\mu F^\mu$的形势。

Noether定理: $$j^\mu = F^\mu - \frac{\partial\lagr}{\partial (\partial_\mu\phi)} \frac{\delta\phi}{\delta\epsilon}$$ 是一个守恒流（即$\partial_\mu j^\mu = 0$)。

\ech
\end{frame}

\begin{frame}
\chtitle{课堂互动}
\bch
试证明Noether定理。
\ech
\end{frame}

\begin{frame}
\chtitle{第一次作业}
\bch
https://github.com/zqhuang/SYSU\_QFTI/homework1.pdf
\ech
\end{frame}

\end{document}
